% -------------------------------------------------------
%  Section 2: Problem statement, motivation, positioning
% -------------------------------------------------------
\section{صورت مسئله، انگیزه و جایگاه در ادبیات}
\label{sec:problem}

مسئله مقاله \textbf{انتساب خانواده بدافزار}\pn{Malware Family Attribution} است، نه تشخیص دودویی بدافزار از نرم‌افزار سالم: ورودی فایل دودویی یک برنامه اجرایی و خروجی توزیع احتمال روی ۹ خانواده است. این تمایز مهم است، چون تمام نمونه‌های \lr{BIG 2015} بدافزار‌اند و هیچ کلاس سالمی در مسئله وجود ندارد؛ پس نتایج مقاله چیزی درباره نرخ هشدار کاذب یک سامانه واقعی نمی‌گویند.

انگیزه روش‌شناختی از دو مشاهده ساخته می‌شود. نخست، به‌دلیل سازوکار هم‌ترازی فایل\pn{File Alignment} در قالب \lr{PE}، انتهای هر بخش با تعدادی بایت صفر پر می‌شود و چون تعداد این بایت‌ها ثابت نیست، مرز میان بخش‌ها در تصویر خاکستری مبهم می‌ماند. دوم، مهاجم می‌تواند ترتیب بخش‌ها را تغییر دهد یا داده زائد تزریق کند و چیدمان فضایی تصویر را که مبنای کار روش‌های پیشین است بر هم بزند.

\lr{Nataraj} و همکاران \cite{nataraj2011images} نخستین بار تبدیل دنباله بایتی به تصویر خاکستری را پیشنهاد کردند و جدول هم‌ترازی عرض بر حسب اندازه فایل از همان کار می‌آید. \lr{Vasan} و همکاران \cite{vasan2020imcfn} تصویر خاکستری را رنگی کردند و یک شبکه پیچشی تنظیم‌شده را روی آن آموزش دادند. نزدیک‌ترین کار، \lr{Xiao} و همکاران \cite{xiao2021section} است که برای هر بخش یک کادر رنگی به تصویر می‌افزایند تا اطلاعات توزیع بخش‌ها برجسته شود؛ ایراد مقاله به این راهکار آن است که کادر رنگی بخشی از پیکسل‌های واقعی را می‌پوشاند. \lr{Chaganti} و همکاران \cite{chaganti2022efficientnet} نیز پیش‌تر نشان داده بودند عرض ثابت سودمند است، ولی دلیل آن را تحلیل نکرده بودند.

جایگاه مقاله در این نقشه روشن است: به‌جای افزودن نشانه بصری به تصویر (راهکار \lr{Xiao})، خود تصویر بر اساس نوع بخش \textbf{تجزیه} می‌شود و زیرتصویر‌ها به کانال‌های مستقل می‌روند؛ یعنی مرزبندی بخش‌ها به‌جای استنتاج از روی تصویر، در ساختار داده ورودی کدگذاری می‌گردد. ادعای «استقلال فضایی» روش نیز از همین‌جا می‌آید. مقاله سه سهم برای خود برمی‌شمارد: صورت‌بندی \lr{Split}/\lr{Mask}، نشان دادن اینکه زیرتصویر‌های بخش‌ها به‌تنهایی نیز ویژگی‌های دسته‌بندی مطلوبی دارند، و طرح عرض تصویر به‌عنوان یک ابرپارامتر مستقل. سهم نخست تازه و به‌روشنی متمایز است. سهم دوم یک مشاهده تجربی ارزشمند است ولی، چنان‌که در \S\ref{sec:critique} و \S\ref{sub:transfer} نشان خواهیم داد، پشتوانه عددی آن شکننده‌تر از چیزی است که متن القا می‌کند. سهم سوم عملاً تأیید مشاهده پیشین \lr{Chaganti} است؛ مقاله به این تقدم اذعان می‌کند، اما تحلیل وعده‌داده‌شده درباره \emph{چرایی} اثر عرض ارائه نمی‌شود و بحث در سطح توصیفی می‌ماند.
